\section{Related work and the evaluation boundary}
\label{sec:related}

\subsection{From task automation to the organization of work}

Task-based accounts separate the displacement of existing work from changes in demand and
the creation of new tasks \citep{acemoglu2019}. That distinction matters for governance:
an automated review can remove execution hours while the organization expands its project
portfolio or assigns the released capacity elsewhere. We study the labor required for a
comparable governed workload, before interpreting actual employment. This fixes an accounting
question rather than assuming that every efficiency gain becomes a redundancy.

Automation can also leave difficult monitoring and intervention work behind
\citep{bainbridge1983}. A routine-case success rate does not say whether the remaining cases
are more expensive than average. The selection identity in Section~\ref{sec:substitution}
makes that concern explicit through the share of baseline labor located in exceptions.
Sequential task assignment and complementary quality requirements motivate treating handoffs
and bottlenecks as part of the process \citep{demirer2026,gans2026}. Our contribution applies
these questions to a defined governance task and reports the quantities needed to move from
an execution claim to a net-work claim. The equations are accounting relationships and
conditional thresholds, not an alternative macroeconomic model.

\subsection{Governance gates and executable process obligations}

Stage-gate organization predates language-model agents \citep{cooper1990}. It supplies a
recognizable decision boundary: a proposal arrives with evidence, reviewers assess it, and
a decision controls the next step. Business-process compliance research studies how obligations
can be represented and monitored during execution \citep{ly2015}. Agentic business-process
management places agents within a broader process-engineering agenda \citep{calvanese2026}.
These traditions provide the process foundations for the task-substitution question studied here.

The question here is what must hold when execution of such a boundary is transferred from
people to software. A review can include both a deterministic policy kernel and interpretation
of evidence whose form is less regular. The kernel may be compiled into ordinary software;
an agent may recover candidate facts, investigate inconsistencies, or assemble a justified
record. The experiment supplies formalized policies and structured facts. Its rules control
identifies the executable decision kernel; the agent runs measure alternative implementations
of the review contract under that same information condition. The contribution of agents to
recovering less structured facts is a separate comparison specified in Section~\ref{sec:implications}.

\subsection{Enterprise-agent benchmarks}

WorkArena evaluates browser-based tasks on enterprise software, using ServiceNow as its
environment \citep{drouin2024}. ITBench concerns operational IT scenarios in site reliability,
compliance and security operations, and financial operations \citep{itbench2025}.
$\tau$-bench evaluates policy-guided interactions among tools, agents, and simulated users,
and explicitly studies consistency across repeated trials \citep{taubench2024}.
These benchmarks establish important neighboring evaluation problems. Table~\ref{tab:related}
locates DGF-Bench by task and acceptance boundary, rather than comparing incompatible scores.

\begin{table}[!htb]
\centering\small
\caption{Neighboring evaluation tasks. The descriptions concern the cited releases; no
cross-benchmark performance ordering is inferred.}
\label{tab:related}
\begin{tabularx}{\linewidth}{@{}P{2.4cm}LL@{}}
\toprule
Benchmark & Central task & Relation to this study\\
\midrule
WorkArena & Execute knowledge-work tasks through enterprise web interfaces & Interface operation can support evidence collection, but is not the gate-review endpoint here\\
ITBench & Address operational IT automation scenarios & Operational execution and remediation differ from reviewing whether a project meets a governance contract\\
$\tau$-bench & Use tools while following domain policies in a user interaction & Motivates repeated reliability and policy compliance; DGF-Bench follows specialist reviews through a dossier route\\
DGF-Bench & Produce a disposition with findings, actions, observed evidence, and authorization & Measures complete review records and whole-route success under explicit supplied policies and facts\\
\bottomrule
\end{tabularx}
\end{table}

The distinction is consequential. A governance agent that asks for a recovery test has
completed one kind of work; an operational agent that repairs the service and demonstrates
recovery has completed another. DGF-Bench scores the former contract and does not take credit
for the latter. Likewise, the current review tools expose structured facts rather than
requiring a browser agent to find every fact in a production interface. These choices define
the measured task; they prevent interpreting a higher percentage as superiority over an
agent evaluated on a broader or less prepared task.

\subsection{Evidence quality and useful agent comparisons}

Work on attribution and citation quality separates a correct answer from a supported one
\citep{gao2023,rashkin2023}. The separation is central to an auditable review: the final
disposition and the record supporting it can fail independently. Our frozen strict endpoint
requires observed sources and conforming excerpts. A separate structural audit tests how
sensitive the conclusions are to field order and numeric representation. It does not convert
lexical or structural matching into an expert judgment about every premise.

Agent evaluation also benefits from simple baselines and explicit resource accounting
\citep{kapoor2024}. We therefore report the deterministic control, failed attempts, inference
costs, and repeated trajectories together with model scores. In this study the strongest
claim is that specified review tasks can execute in software and that their reliability can
be inspected at several levels. Establishing the incremental value of an agent for less
structured evidence requires the matched-information extension in Section~\ref{sec:implications}.
The labor account then supplies a further test: whether the resulting system transfers work
without recreating equivalent hours in verification, exceptions, or maintenance.
